%%%%%%%% ICML 2025 EXAMPLE LATEX SUBMISSION FILE %%%%%%%%%%%%%%%%%

\documentclass{article}

% Recommended, but optional, packages for figures and better typesetting:
\usepackage{microtype}
\usepackage{graphicx}
\usepackage{subfigure}
\usepackage{booktabs} % for professional tables
\usepackage{indentfirst}

% hyperref makes hyperlinks in the resulting PDF.
% If your build breaks (sometimes temporarily if a hyperlink spans a page)
% please comment out the following usepackage line and replace
% \usepackage{icml2025} with \usepackage[nohyperref]{icml2025} above.
\usepackage{hyperref}


% Attempt to make hyperref and algorithmic work together better:
\newcommand{\theHalgorithm}{\arabic{algorithm}}

% Use the following line for the initial blind version submitted for review:
% \usepackage{icml2025}

% If accepted, instead use the following line for the camera-ready submission:
\usepackage[accepted]{icml2025}

% For theorems and such
\usepackage{amsmath}
\usepackage{amssymb}
\usepackage{mathtools}
\usepackage{amsthm}

% if you use cleveref..
\usepackage[capitalize,noabbrev]{cleveref}

%%%%%%%%%%%%%%%%%%%%%%%%%%%%%%%%
% THEOREMS
%%%%%%%%%%%%%%%%%%%%%%%%%%%%%%%%
\theoremstyle{plain}
\newtheorem{theorem}{Theorem}[section]
\newtheorem{proposition}[theorem]{Proposition}
\newtheorem{lemma}[theorem]{Lemma}
\newtheorem{corollary}[theorem]{Corollary}
\theoremstyle{definition}
\newtheorem{definition}[theorem]{Definition}
\newtheorem{assumption}[theorem]{Assumption}
\theoremstyle{remark}
\newtheorem{remark}[theorem]{Remark}

% Todonotes is useful during development; simply uncomment the next line
%    and comment out the line below the next line to turn off comments
%\usepackage[disable,textsize=tiny]{todonotes}
\usepackage[textsize=tiny]{todonotes}


% The \icmltitle you define below is probably too long as a header.
% Therefore, a short form for the running title is supplied here:
\icmltitlerunning{Submission and Formatting Instructions for ICML 2025}

\begin{document}

\twocolumn[
\icmltitle{Modeling Diabetes Risk Using  \\
         Lifestyle and Demographic Predictors}

\begin{icmlauthorlist}
  \icmlauthor{Natalie Ladocsi}{}
  \icmlauthor{Molly O'Leary}{}
  \icmlauthor{Maya Papadopoulos}{}
\end{icmlauthorlist}

% Comment out affiliations (they create the numbered superscripts and footnote)
% \icmlaffiliation{yyy}{Department of XXX, University of YYY, Location, Country}
% \icmlaffiliation{comp}{Company Name, Location, Country}

% Comment out corresponding author footnotes (these show up in the bottom-left)
% \icmlcorrespondingauthor{Firstname1 Lastname1}{first1.last1@xxx.edu}
% \icmlcorrespondingauthor{Firstname2 Lastname2}{first2.last2@www.uk}

% Keywords are fine to keep; they don't print on the page
\icmlkeywords{Machine Learning, ICML}

\vskip 0.3in
]

\setlength{\parindent}{1.5em}  % how big the indent is
\setlength{\parskip}{0pt} 
% this must go after the closing bracket ] following \twocolumn[ ...

% This command actually creates the footnote in the first column
% listing the affiliations and the copyright notice.
% The command takes one argument, which is text to display at the start of the footnote.
% The \icmlEqualContribution command is standard text for equal contribution.
% Remove it (just {}) if you do not need this facility.

%\printAffiliationsAndNotice{}  % leave blank if no need to mention equal contribution
%\printAffiliationsAndNotice{\icmlEqualContribution} % otherwise use the standard text.

\begin{abstract}
The purpose of this project is to estimate the likelihood that an adult develops diabetes in the United States based on demographic and behavioral characteristics. Diabetes is one of the most prevalent chronic diseases in the U.S. and is associated with severe long-term health complications. As rates of obesity and physical inactivity rise, the ability to identify factors that lead to increased risk is critical for prevention and early intervention efforts. Through examining social, behavioral, and biological variables of interest, this project aims to better understand patterns of diabetes risk across the adult population. 

To address this question, we use data from the 2024 National Health Interview Survey (NHIS), which provides information on respondents' health, socioeconomic conditions, and lifestyle behaviors, making it well-suited for analyzing diabetes risk factors. In collecting their data, the NHIS gathers results from around the country. Therefore conclusions drawn reflect larger national trends and are not limited by sampling constraints.

After identifying key variables related to diabetes and its potential predictors, we constructed a unified diabetes indicator variable by combining three separate variables. We merged gestational diabetes, pre-diabetes, and diagnosed diabetes into a single outcome variable. This outcome variable would be used to indicate if a respondent had diabetes. This allowed both diagnosed and at-risk individuals to be captured in the outcome because pre-diabetes and gestational diabetes increases the likelihood of developing type 2 diabetes. 

In order to model diabetes risk, DecisionTreeClassifiers were used due to their ability to handle categorical variables and non-linear relationships. The dataset was split into an 80/20 train-test split that was used universally in all the models to ensure consistency and accuracy in the comparison of results. The two individual decision trees achieved similar accuracy values of around 0.78, but high true positive and low true negative rates indicating instability. Improvement of the model revealed that the optimal maximum tree depth value was 10. In order to address variance found in the confusion matrices of the individual trees, an ensemble of 1,000 trees were bootstrapped to improve stability and slightly increased the accuracy to 0.789. A random forest model reaffirmed these results and identified age and high cholesterol as the most influential predictors. Overall, the models were successful in identifying risk factors, but the precision of individual predictions was limited. 

\end{abstract}

\section{Data Description}
\label{submission}
With the aim of this project being to develop a model that classifies whether or not an adult has diabetes in the United States, the NHIS data was essential to understand. In-person household interviews are conducted and a variety of yes/no, and multiple-choice type questions are used to obtain data. The variables and data collection through this survey on demographic and behavioral characteristics are what we turned to when selecting significant variables. It was first necessary to classify which variables would allow for a determination on if our respondent actually had diabetes. The three variables used to identify a diabetes diagnosis are “Ever had diabetes,” “Ever had prediabetes,” “Ever had gestational diabetes”. For those three data Variables: 
\begin{itemize}
    \item GESDIB\_A: Ever had gestational diabetes with 1 denoting yes and 2 no
    \item PREDIB\_A: Ever had prediabetes with 1 denoting yes and 2 no
    \item DIBEV\_A: Ever had diabetes with 1 denoting yes and 2 no
\end{itemize}
A 1 in the data collection indicates a “Yes”, a 0 indicates a “No”, and 7-9 indicates either “Refused to answer”, “Don’t know”, or “Not ascertained”. For proper cleaning the three variables were combined into one new variable, DIAB\_ANY. Furthermore, we removed all values 7-9, replacing all with the nan value.
Now that our target variable is cleaned and manageable we turned to select our predictor variables. After looking through all the possible variables in the codebook we knew the initial predictor variables must be basic demographic identifiers. We selected age, sex, race, region of household: 
\begin{itemize}
    \item SEX\_A: Sex of sample adult, with 1 denoting a participant coded as male and 2 a participant as female
    \item RACELLP\_A: Race of participant
    \item REGION: Region participant's household is located in
\end{itemize}
Once again we combined 7-9 response options, replacing the values with nan. These factors are a good initial start for a diabetes prediction model. Furthermore, we wanted more specific characteristics for lifestyle habits and situational factors. The factors we chose were, closeness to poverty line, body mass index, education level of sample adults, general health status, frequency of leisurely walks.
\begin{itemize}
    \item RATCAT\_A: Ratio of family income to poverty threshold for patient
    \item EDUCP\_A: Educational level of patient
    \item PHSTAT\_A: General health status
    \item WLKLEIS\_A: If the patient walked for leisure in the past 7 days with 1 denoting yes and 2 no
\end{itemize}
Those four variables were again cleaned to have all the vague 7-9 responses combined and changed to be nan values. 
After a more broad but important look into the lifestyle of the responders through these variables we wanted to target some of the most influential factors that correlate to diabetes. With some basic background research into diabetes we found that cholesterol levels along with levels of hypertension are two of the most determining factors that could predict if an adult could suffer from diabetes in their future. With this in mind we finished off our predictive variable selection with “Ever had high cholesterol?”, “Have you ever been told you had hypertension?”, and the surveyor's BMI and age.  
\begin{itemize}
    \item BMICAT\_A: Categorical body mass index of patient
    \item HYPEV\_A: Ever been told they have hypertension with 1 denoting yes and 2 no
    \item CHLEV\_A: Ever had high cholesterol with 1 denoting yes and 2 no
    \item AGEP\_A: Age of patient
\end{itemize}
These were all continuous variables and therefore did not have to be recoded for vague responses as the other variables above did.
The final predictive factor that needed significant cleaning was regarding the food insecurity variables. We thought determining the connection diabetes had to food insecurity would provide significant information and a more accurate predictive model. Handling these separate variables similarly to the overall diabetes variable, we took 5 variables that all had relevance towards running out of food, worries of being hungry, or used food stamps. Combining these variables we were left with a new variable named, FOOD\_INS. The new code for this variable is “1” indicates food insecurity present, while “0” indicates no food insecurity. One last time the 7-9 response code was replaced with nan values. 

Overall cleaning for this dataset was very straightforward. With our chosen dataset most of the variables are categorical, but given a numerical key. This means the formatting of the binary number keys did not need adjustment because the respondent was not given the freedom to personalize their response. The final step before moving on to coding our model was creating a new table with only the relevant data. The NHIS data is extremely large and therefore we were unable to attach the full CSV to our github repository. By creating a new table with only the relevant data a more manageable file can be analyzed. 
 

\section{Methods and Results}


\subsection{Data Cleaning and Exploratory Analysis}
The dataset was first cleaned by removing columns that were not of interest and keeping only those outlined in the data description above. Our variable of interest, which was the diabetes indicator (DIAB\_ANY) was treated as a binary outcome. If the respondent reported any form of diabetes (gestational, pre-diabetes, or previous diabetes diagnosis) the variable equaled 1. Otherwise its value was 0 indicating that the person was non-diabetic. All variables used to construct the DIAB\_ANY outcome were removed from the predictor set to prevent leakage. Although initially categorical, all of our variables were numerically encoded in the original dataset so those did not need to be re-encoded. After cleaning and subsetting to our variables of interest, the final analytic sample contained 32,629 adults. 

The simple count plot in Figure 1 shows that the majority of respondents do not report diabetes (DIAB\_ANY = 0), while 24.08\% report some form of diabetes. In 2021, crude estimates conducted by the CDC found that 38.4 million people or 11.6\% of the U.S. population had diabetes. This value is lower than that found in our data, but this could be attributed to the fact that those with diabetes are more likely to voluntarily report and fill out the survey than those without. Figure 2 highlights that there is a higher prevalence of diabetes among females (25.08) than males (22.06). In this graph a value of 1 for SEX\_A indicates that the patient is male and 2 indicates that the patient is female. This could also be attributed to the fact that a greater proportion of the patients, with and without diabetes, are female.

\begin{figure}
    \centering
    \includegraphics[width=0.8\linewidth]{Screenshot 2025-12-08 at 9.11.41 PM2.png}
    \caption{Diabetes Prevalence in the Data}
    \label{fig:placeholder}
\end{figure}

\begin{figure}
    \centering
    \includegraphics[width=0.8\linewidth]{Screenshot 2025-12-08 at 9.22.07 PM.png}
    \caption{Diabetes Prevalence by Sex}
    \label{fig:placeholder}
\end{figure}

To investigate variables with continuous value ranges, we examined correlations among age (AGEP\_A), BMI category (BMICAT\_A), self-reported hypertension (HYPEV\_A), and self-reported high cholesterol (CHLEV\_A). These variables show moderate positive associations with diabetes status, aligning with well-known clinical risk factors. Figure 3 which is a correlation heatmap highlights that age and BMI are more strongly associated with DIAB\_ANY, meaning that they are critical to predicting diabetes and were important to include in our tree-based models.

\begin{figure}
    \centering
    \includegraphics[width=0.8\linewidth]{Screenshot 2025-12-08 at 9.24.14 PM.png}
    \caption{Correlation Heatmap of Key Predictors}
    \label{fig:placeholder}
\end{figure}

\subsection{Individual Decision Tree Models}

After initial visualization and exploratory analysis we choose to utilize a decision tree model to further examine our data. A decision tree model was chosen because it handles categorical features, when encoded numerically, extremely efficiently. Furthermore, there is no assumption of linearity that occurs with a decision tree model. While exploring the relationship between diabetes and different demographic characters it was unlikely that the variables affected risk linearly. Additionally, the interpretability of a decision tree model is high, meaning the logic is easily understood. The visualization of this model allows for easy analysis on which variables matter most, the variables closest to the top of the tree are often the most influential. Lastly, utilizing a decision tree model allows analysis into how different variables work together to influence the risk of diabetes. The variables are not all looked at in isolation, but instead the splits reflect combinations of characteristics. 
With that justification in mind we first fit two individual decision trees to predict diabetes status from demographic, health, and food-insecurity-related variables. The cleaned dataset was split into an 80/20 train-test partition and this train-test split was applied to all models and used consistently to ensure accurate comparison. All models were trained on the 80 \% training set and evaluated on the 20 \% test set. The 80/20 split rationale is a standard practice in model learning, but specifically for this data set it was a necessary protocol. Our NHIS dataset is significantly large so this split allows for adequate data for our training and evaluation. 80 \% being our training portion ensures the model has a significant amount of examples to determine the patterns across the different demographic and health characteristics. The 20 \% test set, also allows for statistically reliable estimates as well as detection of overfit. 

Each tree was fit using a DecisionTreeClassifier with:
\begin{itemize}
    \item Max\_depth = 15: to prevent unnecessarily deep trees that could memorize the training data instead of learning the patterns.
    \item Min\_samples\_leaf = 1 \% of training sample size: This prevents the tree from creating trees that are too small or specific, limiting overfitting
\end{itemize}
To illustrate the instability of single trees, we trained two trees on different samples of the training data. However, the same test set was used to evaluate the performance of each tree. Since diabetes status is a binary outcome, the model was evaluated using classification-based metrics rather than regressors like R2. Overall this approach to use multiple trees demonstrates that fundamentally decision tree models have high variance. Therefore for each tree we evaluated:
\begin{itemize}
    \item Accuracy which provides an overall measure of correctness
    \item True Positive Rate (TPR) which describes the proportion of diabetic respondents correctly predicted as diabetic
    \item True Negative Rate (TNR) which describes the proportion of non-diabetic respondents correctly predicted as non-diabetic
    \item Confusion matrices which provide the raw counts that underlie the rates. This allows us to see the absolute number of each error type.
\end{itemize}

\begin{table}
    \centering
    \begin{tabular}{cccl}\toprule
         Model& Accuracy& True Positive Rate & True Negative Rate\\\midrule
         Tree 1&  0.784&  0.932&0.320\\
         Tree 2&  0.782&  0.928&0.322\\ \bottomrule
    \end{tabular}
    \caption{Performance of Single Decision Trees on the 20 \% Test Set}
    \label{tab:placeholder}
\end{table}
Table 1 summarizes the results for the two models. Both trees achieved similar accuracy (~78\%) and high true positive rates (93.2\% and 92.8\%), indicating consistency and strong performance in identifying diabetic patients. This overall shows that our model is finding reliable patterns that separate people with diabetes to those without diabetes. However, the true negative rates differed significantly between the two trees (0.320 vs 0.322), with 0.6\% relative decrease instability in classifying non-diabetic patients was indicated. A slight limitation in the model is revealed that both trees appear heavily biased towards predicting diabetes. 
These conclusions are further justified by the values found in each tree's confusion matrix. To create the confusion matrix, the model’s predictions were compared to the actual diabetes status for each respondent in the 20\% test set. However, the differences in matrix values in Table 2 are extremely small showing that the two individual trees have stability and consistency.  Tree 1 is slightly more sensitive as it identifies 21 more diabetic cases than tree 2. Additionally overall tree 2 is less accurate. 

\begin{table}
    \centering
    \begin{tabular}{cccl}\toprule
         Metric& Tree 1& Tree 2& Change\\\midrule
         True Positive&  4612&  4591& -21\\
         False Negative&  335&  356& +21\\
         False Positive&  1071&  1068& -3\\
         True Negative&  504&  507 & +3\\ \bottomrule
    \end{tabular}
    \caption{Comparison of Confusion Matrices of the Individual Trees}
    \label{tab:placeholder}
\end{table}

Lastly, it is time to dive into the actual decision tree models analysis. It can be determined from looking at the trees in Figures 4 and 5 that they are relatively deep with many lower-level splits. These splits are helpful in capturing the noise in the training data as single trees are high-variance models that are prone to overfitting. The visual patterns within these diagrams are significant. The trees are dominated by orange nodes; this explains why there are high True Positive Rates for both trees, because the orange nodes represent predictors of diabetes, while the blue nodes represent predictions of no diabetes. 

The consistency and stability of the two trees is evident in the similarities of their structure. These similarities are significantly clear in the upper levels of the tree. Another visual factor that appears in this model is the node coloring. The lighter node colors indicate that the node has a mix of both diabetic and non-diabetic people, a more uncertain prediction. This differs from the darker nodes which contain confident predictions in favor or against diabetes risk. 

\begin{figure}
    \centering
    \includegraphics[width=0.9\linewidth]{Screenshot 2025-12-10 at 4.44.54 PM.png}
    \caption{Sample Decision Tree 1}
    \label{fig:placeholder}
\end{figure}

\begin{figure}
    \centering
    \includegraphics[width=0.9\linewidth]{Screenshot 2025-12-10 at 4.45.44 PM.png}
    \caption{Sample Decision Tree 2}
    \label{fig:placeholder}
\end{figure}

Through this model creation there were also some limitations that we came across. As discussed during analysis of Tables 1 and 2 there was an extreme imbalance in the risk determination of diabetes. Although there was almost never a missed case of diabetes this could create false alarms. Finally, there is a potential overfitting in this data model as the structures are complex leading to the assumption that the model didn’t just learn from true patterns but possible noise.

\subsection{Determining Max Depth Value}
In order to validate the choice of tree depth, a test was conducted using the values depth\_values = [5, 10, 15, 20]. Trees with the small depth value of 5, had the lowest test-set accuracy value of 0.784729 which indicates underfitting. Trees with a depth greater than 10 had the same accuracy value as that with a depth of 10 which was 0.786415. Since predictive performance did not improve after a depth of 10, it can be concluded that increasing depth beyond 10 does not improve the model. Therefore a max\_depth of 10 was used in our subsequent models. Table 2 shows the accuracy values for each depth value. It is important to note that there is a very small performance difference between the max\_depth 5 and the rest. Although the depth=10 performs best this small difference indicates that random variation could be the cause, not just true model improvement. 
\begin{table}
    \centering
    \begin{tabular}{cccl}\toprule
         Max\_depth& Accuracy\\\midrule
         5&  0.784729\\
         10&  0.786415\\ 
         15&  0.786415\\ 
         20&  0.786415\\\bottomrule
    \end{tabular}
    \caption{Accuracy Values for Varying Maximum Depth Values}
    \label{tab:placeholder}
\end{table}

\subsection{Bootstrap Aggregation}
To address the instability of single trees, we implemented bootstrap aggregation. This method will construct multiple trees and then average their predictions to reduce overall variance. We constructed an ensemble of T = 1000 decision trees with an upper limit of max\_depth = 10. This is a reduced complexity from the single tree models which had a max\_depth = 15. Each tree in the ensemble was trained on a bootstrap sample of the training data, and predictions of the test set were averaged across trees.
Figure 6 shows a kernel density plot of the distribution of test-set accuracy values across the 1,000 trees that were bootstrapped. The aggregate predicted score across all 1,000 trees or ensemble accuracy has a value 0.788562 and 2.5\% and 97.5\% accuracy quantiles: [0.77736891 0.78702852]. This distribution shows that the ensemble prediction is closer to the upper end of the variability range which is consistent with the theory. 
Compared to the single decision trees the accuracy slightly improved from 0.786 to 0.789. This aligns with the theory that averaging high-variance models reduces variance and improves predictive performance. Stability also greatly increased by taking the average, as the model is no longer relying solely on a particular split chosen by one tree. Overall, bootstrapping improved both predictive accuracy and reliability even though the accuracy value increased only slightly.

\begin{figure}
    \centering
    \includegraphics[width=0.9\linewidth]{Screenshot 2025-12-10 at 6.53.48 PM.png}
    \caption{ Distribution of Test-Set Accuracy Values Across 1,000 Bagged Trees}
    \label{fig:placeholder}
\end{figure}

\subsection{Random Forest and Feature Importance}
To further investigate which predictors drive diabetes risk, we also fit a RandomForestClassifier with max\_depth = 10 on the training set. From our previous model testing we were able to build predictive models towards diabetes risk. We now need to finalize the answer to our initial question and determine the demographic and health characteristics that increase the risk of diabetes. The feature importance allows us to validate our selection of variables and understand the prediction drivers. We computed the accuracy and mean decrease in impurity (MDI) for each feature to visualize these values in the horizontal bar plot in Figure 7.  The accuracy calculated was 0.788 which is similar to that of the bootstrapped ensemble. Figure 7 revealed that:
\begin{itemize}
    \item High-cholesterol (CHLEV\_A) and age (AGEP\_A) are the most influential non-diagnostic predictors of diabetes.
    \item Variables relating to general health measurements and physical activity contributed moderately to predictions.
    \item Sex and food insecurity had the least impact on prediction making but still had some influence which they can not be discounted for.
\end{itemize}

\begin{figure}
    \centering
    \includegraphics[width=0.9\linewidth]{Screenshot 2025-12-10 at 7.00.25 PM.png}
    \caption{Feature Importance Based of MDI Calculations}
    \label{fig:placeholder}
\end{figure}

\subsection{Summary of Findings}
In comparing the models above several patterns, which are consistent with literature knowledge from class emerge. The single decision tree classifiers showed high sensitivity but more variability. From the true positive and true negative calculated values it was determined that individual trees are better at identifying diabetic rather than non-diabetic patients. Analysis of multiple max\_depth values revealed that the optimal tree depth that provides a balance between under and overfitting was 10. Using bootstrapping significantly improved predictive stability and slightly increased accuracy compared to that of single trees. The RandomForestClassifier further confirmed these assumptions while identifying high-cholesterol and age as the dominant predictors of diabetes. 

Confidence in these results is moderate with respect to identifying the most important risk factors, but lower for making individual-level predictions. Increased confidence in the model comes from the knowledge that all three models (single tree, bootstrap, and random forest) produced consistent conclusions regarding high-cholesterol and age. Also, in determining optimal tree depth, it could be seen that the model performance is stable across a range of tree depths. Lastly, bootstrapping reduced variance and further improved accuracy. Overall, the results support high confidence in qualitative conclusions, such as determining risk factors, but lower confidence in quantitative predictions such as determining diabetic outcome. The models are best used for risk-factor identification rather than precise diagnostics.

\section{Conclusion}








\bibliography{example_paper}
\bibliographystyle{icml2025}


%%%%%%%%%%%%%%%%%%%%%%%%%%%%%%%%%%%%%%%%%%%%%%%%%%%%%%%%%%%%%%%%%%%%%%%%%%%%%%%
%%%%%%%%%%%%%%%%%%%%%%%%%%%%%%%%%%%%%%%%%%%%%%%%%%%%%%%%%%%%%%%%%%%%%%%%%%%%%%%
% APPENDIX
%%%%%%%%%%%%%%%%%%%%%%%%%%%%%%%%%%%%%%%%%%%%%%%%%%%%%%%%%%%%%%%%%%%%%%%%%%%%%%%
%%%%%%%%%%%%%%%%%%%%%%%%%%%%%%%%%%%%%%%%%%%%%%%%%%%%%%%%%%%%%%%%%%%%%%%%%%%%%%%
\newpage
\appendix
\onecolumn
\section{You \emph{can} have an appendix here.}

You can have as much text here as you want. The main body must be at most $8$ pages long.
For the final version, one more page can be added.
If you want, you can use an appendix like this one.  

The $\mathtt{\backslash onecolumn}$ command above can be kept in place if you prefer a one-column appendix, or can be removed if you prefer a two-column appendix.  Apart from this possible change, the style (font size, spacing, margins, page numbering, etc.) should be kept the same as the main body.
%%%%%%%%%%%%%%%%%%%%%%%%%%%%%%%%%%%%%%%%%%%%%%%%%%%%%%%%%%%%%%%%%%%%%%%%%%%%%%%
%%%%%%%%%%%%%%%%%%%%%%%%%%%%%%%%%%%%%%%%%%%%%%%%%%%%%%%%%%%%%%%%%%%%%%%%%%%%%%%


\end{document}